\section{Динамика}
%%%%%%%%%%%%%%%%%%%%%%%%%%%%%%%%%%%%%%%%%%%%%%%%%%%%
\subsection{Законы Ньютона}
%%%%%%%%%%%%%%%%%%%%%%%%%%%%%%%%%%%%%%%%%%%%%%%%%%%%
%%%%%%%%%%%%%%%%%%%%%%%%%%%%%%%%%%%%%%%%%%%%%%%%%%%%
\z{Снаряд массой $m = 2\,\text{кг}$ вылетает из ствола орудия в горизонтальном направлении со скоростью $v = 1000\,\text{м/с}$. Определить силу давления $F$ пороховых газов, считая её постоянной, если длина ствола $l = 3.5\,\text{м}$.
}
%
{$F = \frac{mv^2}{2l} = 285\,\text{кН}$}

%%%%%%%%%%%%%%%%%%%%%%%%%%%%%%%%%%%%%%%%%%%%%%%%%%%%
\z{Груз массой $40\,\text{кг}$ поднимают на лёгкой верёвке вертикально вверх так, что скорость груза за $3\,\text{с}$ равномерно возрастает от $0.5\,\text{м/с}$ до $2\,\text{м/с}$. Найти силу натяжения верёвки.
}
%
{$T= m\left(g + \frac{v_2 - v_1}{\Delta t}\right) = 420\,\text{Н}$}

%%%%%%%%%%%%%%%%%%%%%%%%%%%%%%%%%%%%%%%%%%%%%%%%%%%%
\z{Скорость автомобиля изменяется по закону $v(t) = 0.5t$. Найти результирующую силу, действующую на него, если его масса $m = 1\,\text{т}$.
}
%
{$F = 500\,\text{Н}$}

%%%%%%%%%%%%%%%%%%%%%%%%%%%%%%%%%%%%%%%%%%%%%%%%%%%%
\z{Под действием какой горизонтальной силы $F$ вагонетка, масса которой $m = 350\,\text{кг}$, движется по горизонтальным рельсам с ускорением $a = 0.15\,\text{м/с}^2$, если сила сопротивления движению $F_c = 12\,\text{Н}$?
}
%
{$F = ma + F_c = 64.5\,\text{Н}$}

%%%%%%%%%%%%%%%%%%%%%%%%%%%%%%%%%%%%%%%%%%%%%%%%%%%%
\z{Чему равен вес лётчика-космонавта массой $m = 80\,\text{кг}$ при старте ракеты с поверхности Земли вертикально вверх с ускорением $a = 15\,\text{м/с}^2$?
}
%
{$F = 2000\,\text{Н}$}

%%%%%%%%%%%%%%%%%%%%%%%%%%%%%%%%%%%%%%%%%%%%%%%%%%%%
\z{На рисунках представлены силы, действующее на точечное тело массой $m = 200\, \text{г}$ (все силы лежат в одной плоскости, $1 \,\text{клетка} = 1 \,\text{Н}$). С каким ускорением тело движется в каждом случае?
\begin{figure}[!h]
    \centering
    \begin{minipage}[c]{0.32\linewidth} % [c] вместо [t]
        \centering
        \includegraphics[height=2cm]{3.1.6.1} % фиксированная высота
        \caption*{a)}
    \end{minipage}
    \hfill
    \begin{minipage}[c]{0.32\linewidth}
        \centering
        \includegraphics[height=2cm]{3.1.6.2}
        \caption*{b)}
    \end{minipage}
    \hfill
    \begin{minipage}[c]{0.32\linewidth}
        \centering
        \includegraphics[height=2cm]{3.1.6.3}
        \caption*{c)}
    \end{minipage}
\end{figure}
}
%
{\puan $a=20 \,\text{м/с}^2$;~\puan $a=15 \,\text{м/с}^2$;~\puan $a=50 \,\text{м/с}^2$;}

%%%%%%%%%%%%%%%%%%%%%%%%%%%%%%%%%%%%%%%%%%%%%%%%%%%%
\z[{\begin{wrapfigure}[4]{r}{0.3\textwidth}%
		\vspace{-10pt}%
		\begin{flushright}%
 			\includegraphics[width=0.3\textwidth]{3.1.7}%
		\end{flushright}%
	\end{wrapfigure}%
	}]%
{Тело массой $m = 400 \ \text{г}$ движется вдоль оси $Ox$. По приведённому графику зависимости проекции скорости тела $v_x$ от времени $t$ найдите проекцию силы $F_x$, действующей на тело в моменты времени:
	\pu $t_1=2 \, \text{с}$;
	\pu $t_2=4 \, \text{с}$;
	\pu $t_3=6 \, \text{с}$;
	\pu $t_4=8 \, \text{с}$.
	}
%
{\puan $F_1=0\,\text{Н};$~\puan $F_2=2\,\text{Н};$~\puan $F_3=-2\,\text{Н};$~\puan $F_4=-2\,\text{Н}$}
